\section{Auditoría de fuentes y pregunta central: ¿qué amplía realmente la dispersión?}

El título de esta sesión puede inducir con facilidad una impresión errónea: como si bastara con escribir que el modelo tiene «billones» de parámetros para que cada token pase por un cómputo de billones de parámetros. Lo que la sesión discute en realidad es otra cosa: dotar al modelo de un gran \textbf{conjunto de parámetros seleccionables}, pero activar para cada token solo a unos pocos expertos, con lo que se desacopla parcialmente el número total de parámetros del cómputo por token. Entender esto es la condición previa para leer todas las gráficas posteriores de velocidad, calidad y hardware.

Esta reescritura usa la grabación oficial de Stanford Online, los subtítulos oficiales en inglés elaborados manualmente, 38 diapositivas docentes completas recuperadas de la grabación en 1080p y seleccionadas a mano, así como los artículos originales de Switch Transformer, GShard, T5, mT5 y V-MoE. La grabación no publicó un slide PDF independiente, por lo que cada figura indica el tiempo del video; las animaciones repetidas, la portada y las diapositivas a las que se regresó varias veces durante el Q\&A no se insertan de nuevo, pero la explicación oral se conserva en el texto.

\lecturefigure{slide-01-title.jpg}{Tema de la sesión: Scaling Transformers through Sparsity.}{Video oficial de Stanford Online 00:00:06--00:01:41.}

Después de esta diapositiva de título, el expositor no entra directamente en la fórmula de enrutamiento, sino que vuelve primero a los hechos empíricos de dense scaling: la loss de los modelos de lenguaje sigue una tendencia de ley de potencias predecible en función de los parámetros, los datos y el presupuesto de cómputo. La pregunta de investigación de MoE no es negar el dense scaling, sino esta: ¿es posible añadir otra dimensión de escalamiento, la de los «parámetros de uso condicional», manteniendo aproximadamente constante el cómputo por ejemplo?

\lecturefigure{slide-02-neural-scaling-laws.jpg}{Dense neural scaling laws: los modelos más grandes tienen mayor eficiencia de muestra y existe un tamaño óptimo que varía de forma suave con el presupuesto.}{Video oficial de Stanford Online 00:01:42--00:02:35.}

La forma empírica del dense scaling puede expresarse con una ecuación simplificada:

\[
\mathcal{L}(N)=\mathcal{L}_{\infty}+aN^{-\alpha},\qquad a>0,\ \alpha>0.
\]

Aquí, $\mathcal{L}(N)$ denota la pérdida de validación con un tamaño de parámetros $N$; $\mathcal{L}_{\infty}$ es la cota inferior irreducible bajo esa distribución de datos y esos supuestos de modelado; $a$ controla la escala vertical, y $\alpha$ determina la velocidad con la que decae el beneficio marginal de agrandar el modelo. Esta ecuación describe una tendencia empírica; no garantiza que cualquier arquitectura, conjunto de datos o hardware pueda extrapolarse sobre la misma curva.

\teachervoice{El docente enfatiza primero el significado de la scaling law: da a los investigadores la confianza de que «aumentar la escala todavía puede traer beneficios», pero el artículo original estudia principalmente dense model. Por ello, esta sesión presenta la sparsity como un nuevo eje estructural, en lugar de presentar MoE como un módulo local ajeno al scaling.}

\lecturefigure{slide-03-new-axis-sparsity.jpg}{Un nuevo eje de escalamiento: aumentar los parámetros de uso disperso e intentar desacoplarlos del cómputo por ejemplo.}{Video oficial de Stanford Online 00:02:36--00:03:05.}

\begin{importantbox}{El invariante central de la contabilidad de recursos}
Al discutir un sparse model hay que reportar a la vez al menos cinco cantidades: \textbf{número total de parámetros}, \textbf{parámetros activados por token}, \textbf{FLOPs por token}, \textbf{volumen de comunicación entre dispositivos} y \textbf{tiempo wall-clock}. Decir solo «más parámetros con el mismo cómputo» es incompleto, porque los parámetros no activados siguen ocupando el dispositivo o la memoria principal, y el enrutamiento produce además comunicación all-to-all y desbalance de carga.
\end{importantbox}

\begin{warningbox}{Un billón de parámetros no significa que se ejecute un billón de parámetros}
Los 1.6T de Switch-C son el tamaño total del conjunto de parámetros disponibles, no el active path de cada token. A la inversa, que los active FLOPs sean pocos tampoco significa que el sistema no tenga costo: el almacenamiento, la comunicación, los búferes de capacidad, las formas de compilación y la inferencia con baja tasa de procesamiento exponen costos adicionales.
\end{warningbox}

\begin{table}[H]
\centering
\small
\caption{Cantidades de recursos que deben separarse al leer sparse scaling.}
\begin{tabularx}{\textwidth}{L{0.22\textwidth}X X}
\toprule
Cantidad & Pregunta que responde & Lectura errónea frecuente \\
\midrule
Número total de parámetros & ¿Cuántos pesos aprendibles almacena el modelo en total? & Tomarlo como el cómputo por token. \\
Parámetros activados & ¿Por cuáles pesos de expert pasa realmente un token? & Ignorar que tokens distintos siguen rutas distintas. \\
FLOPs & ¿Cuál es la carga aritmética teórica? & Usar los FLOPs como sustituto directo de la latencia real. \\
Volumen de comunicación & ¿Cuántos datos transmiten el token dispatch y la agregación de resultados? & Ignorar que all-to-all se vuelve el cuello de botella. \\
wall-clock & ¿Cuánto tarda en alcanzar la calidad objetivo con un hardware, un lote y una implementación dados? & Generalizar sin considerar el hardware ni el escenario de tasa de procesamiento. \\
\bottomrule
\end{tabularx}
\end{table}

\subsection{Resumen de la sección}

Lo que esta sesión amplía son los \textbf{parámetros disponibles de forma condicional}, no un aumento incondicional del cómputo por token. Todos los resultados posteriores deben interpretarse según estas cinco coordenadas: parámetros totales, parámetros activados, FLOPs, comunicación y tiempo.
